\section{The sharp quaternionic volume bound}
\label{alg:section}

We retain the notation of Section~\ref{geo:section}.  Thus $n=4r$,
$r\ge2$, and
\[
 X=e^2,\qquad L=e^2p_1(E)/4=\alpha X^2+\beta XY,
 \qquad 4\alpha\in\Z,\quad |\beta|=1.
\]
In this section the polynomial degrees of $X$ and $Y$ are both one;
one polynomial degree therefore means cohomological degree four.
Let
\begin{equation}
 \ell(P)=\int_{B} eP(X,Y),
 \qquad \deg P=2r-1.
 \label{alg:functional}
\end{equation}
The inputs of Proposition~\ref{geo:parameters} include
\begin{equation}
 I_j=\ell(X^{2r-1-j}Y^j),\qquad
 \left|\det(I_{i+j})_{0\le i,j<r}\right|=r+1,
 \qquad N=I_0>0.
 \label{alg:middle-input}
\end{equation}
Our goal is
$N\le\binom{4r}{2r}/(2r+1)$, the model value computed in
Proposition~\ref{geo:contact}. The attachment gives polynomial
relations; the middle determinant fixes the scale of the top
functional; an explicit coefficient sum then gives the bound.
We do not prescribe the cohomology ring of the geodesic quotient
in advance.

\subsection{The power relations and the nonzero parameter}

Introduce a formal variable $Z$ with $Z^2=L$ and define homogeneous
polynomials in $X,Y$ by
\begin{align}
 Q(X,Y)&=\frac{(Y+Z)^{r+1}-(Y-Z)^{r+1}}{2Z}
   =\sum_{j\ge0}\binom{r+1}{2j+1}Y^{r-2j}L^j,
   \label{alg:Q}\\
 R(X,Y)&=\frac{(Y+Z)^{r+1}+(Y-Z)^{r+1}}2
   =\sum_{j\ge0}\binom{r+1}{2j}Y^{r+1-2j}L^j.
   \label{alg:R}
\end{align}
Only the terms with valid binomial indices are included.  The quotient
in~\eqref{alg:Q} is a polynomial identity; it involves no division by a
cohomology class.  The degrees of $Q,R$ are $r,r+1$, respectively.

\begin{lemma}
\label{alg:relations}
The classes $Q(X,Y)$ and $R(X,Y)$ vanish in
$H^*(B;\mathbb Q)$.  Moreover $\alpha\ne0$, and consequently
$|4\alpha|\ge1$.
\end{lemma}

\begin{proof}
On the sphere bundle $p:S(E)\to B$, the attachment class is
$x=Y+e\xi$ and satisfies $x^{r+1}=0$.  The fiber antipode sends $\xi$
to $-\xi$, so $(Y-e\xi)^{r+1}=0$ as well.  Taking the sum gives
$p^*R=0$.  The rational Leray--Hirsch decomposition with basis $1,\xi$
makes $p^*$ injective, hence $R=0$.  Taking the difference and integrating
over the sphere fiber, using $p_!\xi=1$, gives $eQ=0$.  The class $Q$ has
cohomological degree $4r=n$, and multiplication by $e$ is injective from
$H^n(B;\mathbb Q)$ to $H^{n+2}(B;\mathbb Q)$ by
Lemma~\ref{geo:gysin}.  Therefore $Q=0$.

Suppose now that $\alpha=0$.  We use the nondegenerate moment pairing
in~\eqref{alg:middle-input} to obtain a contradiction.  Write
\[
 q(y)=Q(1,y),\qquad a(y)=R(1,y),\qquad
 A_0=\C[y]/(q).
\]
The polynomial $q$ has degree $r$ and leading coefficient $r+1$.
The relations $Q=0$ imply that the rule $y^j\mapsto I_j$, for
$0\le j\le2r-1$, agrees with reduction modulo $q$ followed by a
functional $F:A_0\to\C$.  Indeed, every multiple of $q$ needed in
dividing a polynomial of this degree has multiplier of degree at most
$r-1$, and its homogenization is a top-degree multiple of $Q$.
Thus the matrix of $(f,g)\mapsto F(fg)$ in $1,y,\ldots,y^{r-1}$ is
precisely the nonsingular matrix in~\eqref{alg:middle-input}.
The relation $R=0$ similarly gives
\begin{equation}
 F(a(y)y^j)=0\qquad(0\le j\le r-2).
 \label{alg:orthogonal-R}
\end{equation}

Let $\varepsilon:A_0\to\C$ extract the coefficient of $y^{r-1}$
in the remainder of degree less than $r$.  Its product pairing is
nondegenerate: the entry $\varepsilon(y^{i+j})$ is zero when
$i+j<r-1$ and one when $i+j=r-1$, for $0\le i,j<r$.
Reversing the columns gives a triangular matrix with diagonal entries one.
Hence there is a unique $b\in A_0$ such that
$F(f)=\varepsilon(bf)$ for all $f$.  Nondegeneracy of the $F$ pairing
means that multiplication by $b$ is invertible, so $b$ is a unit.
For the $\varepsilon$ pairing the orthogonal complement of
$1,y,\ldots,y^{r-2}$ consists exactly of the constants.  To see this,
write a remainder $c_0+\cdots+c_{r-1}y^{r-1}$ and successively test
against $1,y,\ldots,y^{r-2}$; the tests force
$c_{r-1},c_{r-2},\ldots,c_1$ to vanish in that order.
Equation~\eqref{alg:orthogonal-R} therefore says that $ba$ is constant.
In particular any nonzero class $a$ satisfying that equation must be a
unit: its product with the unit $b$ is a nonzero constant.

But the class of $a=R(1,y)$ is nonzero and is not a unit.  Both $q$ and
$a$ vanish at $y=0$ when $\alpha=0$, so evaluation at zero precludes
invertibility.  To check that the class is nonzero, choose an
$(r+1)$st root of unity $\eta\ne1,-1$, and put
\[
 c_\eta=\frac{\eta-1}{\eta+1},\qquad
 y_0=\frac{\beta}{c_\eta^2},\qquad Z_0=c_\eta y_0.
\]
These numbers are nonzero, $Z_0^2=\beta y_0$, and
$y_0+Z_0=\eta(y_0-Z_0)$ with $y_0-Z_0\ne0$.  The formulas for $Q,R$
give $q(y_0)=0$ but $a(y_0)=(y_0-Z_0)^{r+1}\ne0$.  Thus $a$ cannot
be zero modulo $q$, a contradiction.  We conclude that $\alpha\ne0$.
Its stated integrality then gives $|4\alpha|\ge1$.
\end{proof}

\begin{lemma}
\label{alg:complete-intersection}
For $\alpha\ne0$, the polynomials $Q,R$ have no common nonzero zero
over $\C$.  The algebra
\[
 \mathcal A=\C[X,Y]/(Q,R)
\]
has Hilbert series
\begin{equation}
 \frac{(1-t^r)(1-t^{r+1})}{(1-t)^2}
   =(1+t+\cdots+t^{r-1})(1+t+\cdots+t^r).
 \label{alg:hilbert}
\end{equation}
In particular its top degree is $2r-1$, with dimension one, and
\begin{equation}
 X:\mathcal A_{r-1}\longrightarrow\mathcal A_r
 \quad\hbox{is an isomorphism},\qquad
 \dim\mathcal A_{r-1}=\dim\mathcal A_r=r.
 \label{alg:middle-isomorphism}
\end{equation}
Furthermore the complete rational relation ideal of the classes
$e,Y$ is generated by $Q(e^2,Y)$ and $R(e^2,Y)$.
\end{lemma}

\begin{proof}
At a common zero $(X,Y)$ choose a complex square root $Z$ of
$\alpha X^2+\beta XY$.  If $Z\ne0$, simultaneous vanishing of $Q,R$
would force both $(Y+Z)^{r+1}$ and $(Y-Z)^{r+1}$ to vanish, which is
impossible unless $Y=Z=0$.  If $Z=0$, the formulas become
$Q=(r+1)Y^r$ and $R=Y^{r+1}$, so $Y=0$; then $\alpha X^2=0$
forces $X=0$ as well.

The two homogeneous polynomials are therefore relatively prime.
Multiplication by $Q$ in $\C[X,Y]$ is injective, and multiplication by
$R$ is injective modulo $Q$: in the polynomial unique factorization
domain, $Q\mid Rf$ and relative primality imply $Q\mid f$.
The two resulting multiplication exact sequences give
\eqref{alg:hilbert}.  On setting $X=0$, the quotient becomes
$\C[Y]/(Y^r)$.  Thus the degree-$r$ cokernel of multiplication by $X$
vanishes.  The two dimensions read from~\eqref{alg:hilbert} are equal,
which proves~\eqref{alg:middle-isomorphism}.

For the assertion about the whole cohomology ring,
Lemma~\ref{geo:gysin} gives generation of $H^*(B;\mathbb Q)$
by $e,v$ and total dimension $2r(r+1)$.
Since $Y=v+(b/2)e^2$ by \eqref{geo:Y-lift}, the same ring is
generated by $e,Y$.

The polynomial algebra $\C[e,Y]$ is free on $1,e$ over $\C[X,Y]$
under $X=e^2$.  Quotienting by the two displayed relations therefore
gives $\mathcal A\oplus e\mathcal A$, of dimension $2r(r+1)$ by
\eqref{alg:hilbert}.  Its surjection to $H^*(B;\C)$ is an
isomorphism by the equality of dimensions.  The relations and maps are
defined over $\mathbb Q$, giving the rational assertion as well.
\end{proof}

\subsection{An explicit top functional and its determinant}

The algebraic method is the classical perfect-pairing duality for
finite graded complete intersections; see
\citet[Corollary~3.5 and the discussion following Theorem~9.1]
{Huneke1999}.
We make the pairing explicit to determine its normalization and
the required determinant.

The power relations suggest replacing $Y+Z$ and $Y-Z$ by two nilpotent
variables $a_1,a_2$.  With $Y=(a_1+a_2)/2$ and $Z=(a_1-a_2)/2$,
the equation $Z^2=L$ becomes a quadratic relation for $X$.
The resulting algebra admits explicit coefficient extraction; the
involution exchanging $a_1,a_2$ will recover the pairing on $\mathcal A$
without cancelling $Z$.

For any complex $\alpha\ne0$ and any complex $\beta$, set
\begin{equation}
 \mathcal D_{\alpha,\beta}=
 \frac{\C[a_1,a_2,X]}
 {\bigl(a_1^{r+1},a_2^{r+1},
  4\alpha X^2+2\beta(a_1+a_2)X-(a_1-a_2)^2\bigr)}.
 \label{alg:extraction-algebra}
\end{equation}
All three variables have polynomial degree one.  Every class has a
unique remainder that has $X$-degree at most one and each $a_i$-degree
at most $r$.  Denote extraction of $Xa_1^ra_2^r$ in that remainder by
$\varepsilon_{\mathcal D}$.  On homogeneous polynomials of degree
$2r-1$ define
\begin{equation}
 \Phi_{\alpha,\beta}(P)=
 \varepsilon_{\mathcal D}\left(
   (a_1-a_2)^2P\left(X,\frac{a_1+a_2}{2}\right)\right).
 \label{alg:Phi}
\end{equation}

\begin{lemma}
\label{alg:determinant}
The functional $\Phi_{\alpha,\beta}$ descends to the top degree of
$\mathcal A$ and is nonzero.  If
\[
 J_j=\Phi_{\alpha,\beta}(X^{2r-1-j}Y^j),
\]
then for every $\alpha\ne0$ and every $\beta\in\C$,
\begin{equation}
 \det(J_{i+j})_{0\le i,j<r}
 =(-1)^{r(r+1)/2}(r+1)(4\alpha)^{-r(r-1)/2}.
 \label{alg:determinant-formula}
\end{equation}
For the geometric parameters there is consequently a scalar $K$ with
\begin{equation}
 \ell=K\Phi_{\alpha,\beta},\qquad
 |K|=|4\alpha|^{(r-1)/2}.
 \label{alg:normalization}
\end{equation}
\end{lemma}

\begin{proof}
Put $Y=(a_1+a_2)/2$ and $Z=(a_1-a_2)/2$ inside $\mathcal D$.
Its quadratic relation becomes $Z^2=L$ and its other relations become
\[
 R=0,\qquad ZQ=0.
\]
It follows at once that multiplying either $R$ or $Q$ by $4Z^2$ gives
zero, so~\eqref{alg:Phi} annihilates the top-degree multiples of both
relations.

We next prove nondegeneracy, uniformly in the complex parameter $\beta$.
The coefficient functional on
\[
 D_0=\C[a_1,a_2]/(a_1^{r+1},a_2^{r+1})
\]
extracting $a_1^ra_2^r$ has a perfect product pairing, since every
monomial pairs with its complementary monomial.  The quadratic relation
in~\eqref{alg:extraction-algebra} is monic after division by $4\alpha$,
so $\mathcal D=D_0\oplus XD_0$.  Its coefficient pairing is perfect as
well.  Indeed, if $f+Xg$ annihilates every element of $D_0$, extraction
shows that $g=0$; testing then against $XD_0$ shows that $f=0$.
The grading makes this a perfect pairing between complementary degrees
summing to $2r+1$.

The involution exchanging $a_1,a_2$ preserves
$\varepsilon_{\mathcal D}$.  Its even and odd subspaces are therefore
orthogonal, and the odd subspaces in complementary degrees pair
perfectly.  Their exact module description is particularly useful.
Before imposing $R,ZQ$, the relation $Z^2=L$ gives the direct sum
\[
 \C[X,Y]\oplus Z\C[X,Y].
\]
The ideal generated by $R,ZQ$ has even part $(R,LQ)$ and odd part
$Z(R,Q)$.  Consequently
\begin{equation}
 \mathcal D^-\simeq Z\mathcal A
 \quad\hbox{as a graded module, with a shift by one}.
 \label{alg:odd-module}
\end{equation}
This follows by taking the two summands of the ideal; no cancellation
of $Z$ in the quotient is used.

By~\eqref{alg:middle-isomorphism}, multiplication by $X$ carries
$\mathcal D^-_r$ isomorphically onto $\mathcal D^-_{r+1}$.
In the basis
\[
 (a_1-a_2)X^{r-1-i}Y^i,\qquad 0\le i<r,
\]
of $\mathcal D^-_r$, the bilinear pairing
$(f,g)\mapsto\varepsilon_{\mathcal D}(fXg)$ has matrix
$(J_{i+j})$.  It is nondegenerate because the odd complementary-degree
pairing is perfect.  This proves in particular that $\Phi$ is nonzero.

For fixed nonzero $\alpha$, reduction of powers of $X$ by its monic
quadratic relation makes every $J_j$ a polynomial in $\beta$.
The determinant is thus a polynomial in $\beta$ which, by the preceding
argument, has no complex zero.  It is a nonzero constant in $\beta$.
It remains to evaluate it at $\beta=0$.

First set $\alpha=1/4$ as well, so that $X^2=(a_1-a_2)^2$.
In $\mathcal A$ use the linear coordinates
\[
 A_1=Y+X/2,\qquad A_2=Y-X/2.
\]
For $p+q=2r-1$, the coefficient of $X$ in the quadratic remainder of
$A_1^pA_2^q$ is the polynomial
\[
 \frac{a_1^pa_2^q-a_1^qa_2^p}{2(a_1-a_2)}.
\]
The numerator is divisible by $a_1-a_2$.  Multiplying by
$(a_1-a_2)^2$ and extracting $a_1^ra_2^r$ gives
\begin{equation}
 \Phi_{1/4,0}(A_1^pA_2^q)=
 \begin{cases}
 -1,&(p,q)=(r,r-1),\\
 1,&(p,q)=(r-1,r),\\
 0,&\text{otherwise}.
 \end{cases}
 \label{alg:base-functional}
\end{equation}
In the degree-$(r-1)$ monomial basis
$A_1^{r-1-i}A_2^i$, the matrix of the middle $X=A_1-A_2$ pairing
has entries $-2$ for $i+j=r-1$, entries $1$ for $i+j=r-2$ or $r$,
and zero otherwise.  Reversing columns gives the tridiagonal matrix
with diagonal $-2$ and adjacent entries $1$.
Its size-$r$ determinant is $(-1)^r(r+1)$: expansion along the last
row gives the recurrence $d_r=-2d_{r-1}-d_{r-2}$ with $d_0=1,d_1=-2$.
Column reversal has sign $(-1)^{r(r-1)/2}$.
The linear change from $X,Y$ to $A_1,A_2$ has determinant one, so its
degree-$(r-1)$ monomial change also has determinant one.  We have proved
\eqref{alg:determinant-formula} at $(\alpha,\beta)=(1/4,0)$.

At $\beta=0$ for general $\alpha\ne0$, odd-indexed $J_j$ vanish.
For $j=2h$, quadratic reduction gives
\[
 J_{2h}(\alpha,0)
   =(4\alpha)^{-(r-1-h)}J_{2h}(1/4,0).
\]
Every nonzero permutation term in the size-$r$ determinant therefore
has parameter exponent
\[
 -\sum_{i=0}^{r-1}\left(r-1-\frac{i+\sigma(i)}2\right)
   =-\frac{r(r-1)}2.
\]
This proves~\eqref{alg:determinant-formula} for $\beta=0$, and its
independence of $\beta$ proves the full identity.

Finally, $\ell$ and $\Phi$ are nonzero functionals on the same
one-dimensional top degree of $\mathcal A$, so $\ell=K\Phi$.
Taking absolute determinants in~\eqref{alg:middle-input} and
\eqref{alg:determinant-formula} gives
\[
 r+1=|K|^r(r+1)|4\alpha|^{-r(r-1)/2},
\]
which is~\eqref{alg:normalization}.
\end{proof}

\subsection{The model upper bound}

\begin{theorem}
\label{alg:upper}
For a quaternionic-type Blaschke manifold of dimension $4r$, $r\ge2$,
the contact characteristic number of Proposition~\ref{geo:contact}
satisfies
\begin{equation}
 N=\int_{B}e^{4r-1}
 \le \frac1{2r+1}\binom{4r}{2r}.
 \label{alg:sharp-upper}
\end{equation}
\end{theorem}

\begin{proof}
The quadratic relation in~\eqref{alg:extraction-algebra} is
\[
 X^2=\mathfrak a X+\mathfrak b,\qquad
 \mathfrak a=-\frac{\beta}{2\alpha}(a_1+a_2),\quad
 \mathfrak b=\frac{(a_1-a_2)^2}{4\alpha}.
\]
The coefficient of $X$ in the remainder of $X^q$ satisfies
$D_1=1$, $D_2=\mathfrak a$, and $D_q=\mathfrak a D_{q-1}+\mathfrak b D_{q-2}$.
Induction on this recurrence yields
\[
 D_{2r-1}=
 \sum_{j=0}^{r-1}\binom{2r-2-j}{j}\mathfrak a^{2r-2-2j}\mathfrak b^j.
\]
Consequently, with ordinary polynomial coefficient extraction,
\begin{equation}
 \Phi_{\alpha,\beta}(X^{2r-1})
 =\sum_{j=0}^{r-1}\binom{2r-2-j}{j}
   \left(-\frac{\beta}{2\alpha}\right)^{2r-2-2j}
   (4\alpha)^{-j}c_j,
 \label{alg:finite-sum}
\end{equation}
where
\[
 c_j=[a_1^ra_2^r]
       (a_1+a_2)^{2r-2-2j}(a_1-a_2)^{2j+2}.
\]
The exact sign of this central coefficient is
\begin{equation}
 \operatorname{sign}(c_j)=(-1)^{j+1}.
 \label{alg:central-sign}
\end{equation}
Indeed, substitute $a_1=e^{it}$, $a_2=e^{-it}$ and average over a full
period.  The coefficient becomes $(-1)^{j+1}$ times the strictly positive
integral of a positive constant times
$\cos^{2r-2-2j}(t)\sin^{2j+2}(t)$.

By Lemmas~\ref{alg:relations} and~\ref{alg:determinant}, the parameter
factor in the absolute value of each summand of
$K\Phi_{\alpha,\beta}(X^{2r-1})$, apart from its binomial and central
coefficients, is
\begin{align*}
 |K|\left|\frac{\beta}{2\alpha}\right|^{2r-2-2j}
       |4\alpha|^{-j}
 &=2^{2r-2-2j}|4\alpha|^{j-3(r-1)/2}\\
 &\le 2^{2r-2-2j}.
\end{align*}
Here $|\beta|=1$, $|4\alpha|\ge1$, and
$j-3(r-1)/2\le-(r-1)/2<0$.
The bound on the right is attained by the formal parameters
$\alpha=-1/4$, $\beta=1$, for which the normalization factor
$|4\alpha|^{(r-1)/2}$ is one.
At those parameters~\eqref{alg:central-sign} shows that all the summands
in~\eqref{alg:finite-sum} have the same, negative, sign.  The sum of
the termwise absolute bounds therefore equals the absolute value of
that single formal evaluation.

We evaluate it explicitly.  When $\alpha=-1/4$, $\beta=1$, the two
formal roots of the quadratic in $X$ are
$(\sqrt{a_1}+\sqrt{a_2})^2$ and
$(\sqrt{a_1}-\sqrt{a_2})^2$.  Their divided difference gives the
polynomial
\begin{align*}
 D_{2r-1}
 &=\frac{(\sqrt{a_1}+\sqrt{a_2})^{4r-2}
          -(\sqrt{a_1}-\sqrt{a_2})^{4r-2}}
         {4\sqrt{a_1a_2}}\\
 &=\frac12\sum_{h=0}^{2r-2}\binom{4r-2}{2h+1}
                 a_1^{2r-2-h}a_2^h.
\end{align*}
Thus multiplication by $(a_1-a_2)^2$ and central coefficient
extraction give
\[
 \left|\Phi_{-1/4,1}(X^{2r-1})\right|
 =\binom{4r-2}{2r-1}-\binom{4r-2}{2r-3}
 =\frac1{2r+1}\binom{4r}{2r}.
\]
Since $N=\ell(X^{2r-1})>0$, the preceding termwise comparison proves
\eqref{alg:sharp-upper}.
\end{proof}
